\documentclass[UTF8,a4paper,11pt,openany]{ctexbook}
\usepackage{../common/statlearnbook}
\SLSetBookMetadata{第 5 册}{统计学习方法讲义 v2.6.0}{采样方法、主题模型与图排序}
\begin{document}
\setcounter{page}{830}
\setcounter{chapter}{34}
\setcounter{figure}{3}
\setcounter{equation}{3}
\pagestyle{headings}
\chaptermark{狄利克雷分布与共轭先验}
\sectionmark{狄利克雷分布}
图\ref{fig:V5-C05-simplex-geometry}说明三分类概率向量虽然写成三个分量，却只有两个自由度。

\paragraph{几何面：概率向量住在哪里？}
% v2.3.1 figure UID: FIG-P660-01
% Proposed post-v2.3.1 polish for FIG-P660-01: protect key-label size and stroke hierarchy.
\tikzset{
  slfig-FIG-P660-01/.style={font=\fontsize{9.2pt}{11.0pt}\selectfont,every path/.append style={line cap=round,line join=round}},
  slfig-FIG-P660-01-boundary/.style={draw=SLMidBlue,line width=0.90pt},
  slfig-FIG-P660-01-grid/.style={draw=SLRuleGray,line width=0.40pt},
  slfig-FIG-P660-01-component/.style={draw=SLTeal,line width=0.58pt,densely dashed},
  slfig-FIG-P660-01-card/.style={rounded corners=2pt,text width=52mm,align=left,inner xsep=3.4mm,inner ysep=2.1mm}
}
\pgfkeysifdefined{/tikz/alt}{}{\tikzset{alt/.code={}}}
\begin{figure}[htbp]
\centering
\begin{tikzpicture}[slfig-FIG-P660-01,>=Stealth,every node/.style={font=\fontsize{9.5pt}{11.4pt}\selectfont},
  >={Stealth[length=2.3mm,width=1.6mm]},
  every path/.append style={line width=0.55pt},
  every node/.append style={line width=0.8pt},
  alt={对象—关系—结论：三类别概率向量位于二维单纯形上，三个顶点分别代表一个类别概率为1；内部点的重心坐标就是三个类别概率，因此Dirichlet分布虽写在三维坐标中，实际支撑集只有二维}]
\coordinate (e1) at (0,0);
\coordinate (e2) at (7,0);
\coordinate (e3) at (3.5,6.062);
\coordinate (th) at (3.85,3.031);

\fill[SLSoftBlue] (e1)--(e2)--(e3)--cycle;
\draw[slfig-FIG-P660-01-boundary] (e1)--(e2)--(e3)--cycle;
\foreach \t in {0.2,0.4,0.6,0.8}{
  \draw[slfig-FIG-P660-01-grid]
    ($(e1)!\t!(e2)$) -- ($(e3)!\t!(e2)$);
  \draw[slfig-FIG-P660-01-grid]
    ($(e2)!\t!(e3)$) -- ($(e1)!\t!(e3)$);
  \draw[slfig-FIG-P660-01-grid]
    ($(e3)!\t!(e1)$) -- ($(e2)!\t!(e1)$);
}
\coordinate (q1) at ($(e2)!(th)!(e3)$);
\coordinate (q2) at ($(e1)!(th)!(e3)$);
\coordinate (q3) at ($(e1)!(th)!(e2)$);
\draw[slfig-FIG-P660-01-component] (th)--(q1)
  node[anchor=west,inner sep=0pt,font=\fontsize{8.7pt}{10.4pt}\selectfont] at (5.15,4.08) {$\theta_1=.2$};
\draw[slfig-FIG-P660-01-component] (th)--(q2)
  node[anchor=east,inner sep=0pt,font=\fontsize{8.7pt}{10.4pt}\selectfont] at (1.95,4.08) {$\theta_2=.3$};
\draw[slfig-FIG-P660-01-component] (th)--(q3)
  node[anchor=north west,inner sep=0pt,font=\fontsize{8.7pt}{10.4pt}\selectfont] at (4.15,-0.32) {$\theta_3=.5$};
\fill[SLDeepBlue] (th) circle (3.2pt);

\node[below=1.6mm,align=center] at (e1) {$e_1=(1,0,0)$\\类别1确定};
\node[below=1.6mm,align=center] at (e2) {$e_2=(0,1,0)$\\类别2确定};
\node[above=1.6mm,align=center] at (e3) {$e_3=(0,0,1)$\\类别3确定};
\node[anchor=west,inner sep=0pt] at (5.55,3.10)
  {$\theta=(0.2,0.3,0.5)$};
\node[slfig-FIG-P660-01-card,draw=SLRuleGray,fill=white] at (12.25,4.45)
  {$\Delta^2=\{\theta\in\mathbb R^3:\theta_i\ge0,\ \sum_i\theta_i=1\}$\\
   $\dim(\Delta^2)=2$};
\node[slfig-FIG-P660-01-card,draw=SLRuleGray,fill=SLSoftGray] at (12.25,2.20)
  {内部：三个分量均为正；\\边：一个分量为零；\\顶点：一个类别概率为1。};
\node[slfig-FIG-P660-01-card,draw=SLMidBlue,fill=SLSoftBlue] at (12.25,-0.05)
  {三元概率向量有三个坐标，却只具有两个自由度；三角形中的位置就是重心坐标。};
\end{tikzpicture}
\caption{三类别概率向量位于二维单纯形上，三个顶点分别代表一个类别概率为1；内部点的重心坐标就是三个类别概率，因此Dirichlet分布虽写在三维坐标中，实际支撑集只有二维}
\label{fig:V5-C05-simplex-geometry}
\end{figure}

\FloatBarrier

\begin{theorem}[Gamma归一化构造]
设$G_i$彼此独立且$G_i\sim\operatorname{Gamma}(\alpha_i,b)$。令$S=\sum_iG_i$、$\Theta_i=G_i/S$，则$\boldsymbol\Theta\sim\operatorname{Dir}(\boldsymbol\alpha)$，并且$\boldsymbol\Theta$与$S$独立。
\end{theorem}
\end{document}
